\FloatBarrier

\section{HTCC CAD cone transfer and validation}
\label{sec:htcc}

The October 19 HTCC update transfers three cured GEMC3 meshes into the GEMC2 clas12Tags repository. These are
the HTCC entry cone, Moller cone, and Moller-cone extension; they are separate from the beamline Moller-cone
studies above. The active source and runtime copies match the GEMC3 assets byte for byte, and the unmodified
GEMC2 engineering inputs are archived separately. Names, Rohacell material, passive 180-degree y rotations,
and all run-dependent placements remain unchanged. The variations default, rga\_spring2018, and rga\_fall2018
use runs 11, 3029, and 4763, respectively; the extension retains its additional 1 mm z shift.

Curing merges vertices, closes mesh boundaries, and repairs orientation. The entry cone is simplified to 6,000
triangles; the Moller cone additionally uses self-intersection removal and finishes at 5,866 triangles. The
already-light extension retains 1,748 triangles. No additional engineering-hole filling or Poisson
reconstruction is applied. Total stored triangles fall from 26,798 to 13,614, a 49.20\% reduction.
Table~\ref{tab:htcc-mesh} lists the topology counts from the same PyMeshLab inspection.


\begin{table}[htbp]
\centering
\caption{HTCC cone topology before and after the transferred curing and simplification.}
\label{tab:htcc-mesh}
\setlength{\tabcolsep}{5pt}
\begin{tabular}{llrrr}
\toprule
Mesh & Input & Triangles & Boundary edges & Non-manifold edges \\
\midrule
\texttt{htccCone} & Raw & 18,818 & 144 & 0 \\
\texttt{htccCone} & Cured & 6,000 & 0 & 0 \\
\texttt{htccMollerCone} & Raw & 6,232 & 0 & 45 \\
\texttt{htccMollerCone} & Cured & 5,866 & 0 & 24 \\
\texttt{htccMollerConeExt} & Raw & 1,748 & 0 & 0 \\
\texttt{htccMollerConeExt} & Cured & 1,748 & 0 & 0 \\
\bottomrule
\end{tabular}
\end{table}

All three cured meshes have zero boundary edges. The entry cone and extension have no non-manifold edges,
while the Moller cone retains 24, reduced from 45 in the raw input. These edge counts depend on the mesh
import and vertex-merging convention and should not be equated to affected-face counts. The largest change
among the six axis-aligned entry-cone bounds is 0.737915 mm. The Moller cone and extension retain their
original bounds exactly. Bounds agreement and zero boundary edges do not establish material-volume or surface
equivalence, and the remaining non-manifold Moller-cone edges are an unresolved topology limitation.

Validation uses GEMC 6.2 with Geant4 10.7.4, a local CCDB SQLite database, no magnetic field, no GUI, no event
output, and zero transported events. Native-plus-CAD initialization uses STD + FTFP\_BERT + Optical. All three
variations complete through both TEXT/CAD and SQLite/CAD, giving six successful initializations. For each
variation, the 537 native TEXT rows match the effective SQLite geometry rows field by field; the three GXML
CAD placements match their effective SQLite CAD rows. Native geometry and materials, CAD placements, and the
SQLite database were retained by the transfer.

The full default SQLite overlap check reports entry-cone contact with htccEntryWindow. An isolated cone/window
fixture uses STD + FTFP\_BERT, CHECK\_OVERLAPS=1000, RANDOM=12345, and the unchanged window geometry and cone
placement. The cured cone produces one GeomVol1002 exception block, with reported penetration of at least
52.4771 micrometres; the matched raw-cone check produces none. The full default check reports at least 52.6119
micrometres for the same pair, approximately 0.053 mm. The two Moller cones report no overlap in that default
assembly check. The entry-window overlap remains unresolved in the transferred mesh; successful construction
does not certify an overlap-free assembly. Sampled penetration reports are not exact overlap depths or
volumes.

The retained zero-event logs verify construction and provide geometry diagnostics. They are not repeated
loading or tracking benchmarks, so no CPU, wall-time, memory, or transport-speed gain is assigned to the
49.20\% triangle reduction. Photon yield, material budget, detector response, and particle navigation
equivalence require separate studies. The HTCC archive retains raw and cured STLs, input fixtures, scripts,
topology hashes, validation logs, and database-row snapshots.
